% !TeX TS-program = pdflatex

\documentclass[a4paper]{article}

% \usepackage[default]{fontsetup}

\usepackage{fancyhdr}
\usepackage{extramarks}
\usepackage{amsmath}
\usepackage{amsthm}
\usepackage{amsfonts}
\usepackage{tikz}
\usepackage[plain]{algorithm}
\usepackage{algpseudocode}
\usepackage{enumerate}
\usepackage{tikz}

\usetikzlibrary{automata,positioning}

%
% Basic Document Settings
%  

\topmargin=-0.2in
\evensidemargin=0in
\oddsidemargin=0in
\textwidth=6.5in
\textheight=9.5in
\headsep=0.25in

\linespread{1.1}

\pagestyle{fancy}
\lhead{\hmwkAuthorName}
\chead{\hmwkClass : \hmwkTitle}
\rhead{\firstxmark}
\lfoot{\lastxmark}
\cfoot{\thepage}

\renewcommand\headrulewidth{0.4pt}
\renewcommand\footrulewidth{0.4pt}

\setlength\parindent{0pt}

%
% Create Problem Sections
%

%修改下标的地方
%\newcommand{\enterProblemHeader}[1]{
	%	\nobreak\extramarks{}{Problem \arabic{#1} continued on next page\ldots}\nobreak{}
	%	\nobreak\extramarks{Problem \arabic{#1} (continued)}{Problem \arabic{#1} continued on next page\ldots}\nobreak{}
	%}

\newcommand{\exitProblemHeader}[1]{
    \nobreak\extramarks{Problem \arabic{#1} }{Problem \arabic{#1} continued on next page\ldots}\nobreak{}
	\stepcounter{#1}
	\nobreak\extramarks{Problem \arabic{#1}}{}\nobreak{}
}

\newcommand*\circled[1]{\tikz[baseline=(char.base)]{
		\node[shape=circle,draw,inner sep=2pt] (char) {#1};}}


\setcounter{secnumdepth}{0}
\newcounter{partCounter}
\newcounter{homeworkProblemCounter}
\setcounter{homeworkProblemCounter}{1}
\nobreak\extramarks{Problem \arabic{homeworkProblemCounter}}{}\nobreak{}

%
% Homework Problem Environment
%
% This environment takes an optional argument. When given, it will adjust the
% problem counter. This is useful for when the problems given for your
% assignment aren't sequential. See the last 3 problems of this template for an
% example.
%

\newenvironment{homeworkProblem}[1][-1]{
	\ifnum#1>0
	\setcounter{homeworkProblemCounter}{#1}
	\fi
	\section{Problem \arabic{homeworkProblemCounter} }
	\setcounter{partCounter}{1}
	% \enterProblemHeader{homeworkProblemCounter} % 注释掉未定义的命令
}{
	\exitProblemHeader{homeworkProblemCounter}
}

%
% Homework Details
%   - Title
%   - Class
%   - Due date
%   - Name
%   - Student ID

\newcommand{\hmwkTitle}{Homework\ \#04}
\newcommand{\hmwkClass}{Probability \& Statistics for EECS}
\newcommand{\hmwkDueDate}{2026-10-20}
\newcommand{\hmwkAuthorName}{}
\newcommand{\hmwkAuthorID}{}


%
% Title Page
%

\title{
	\vspace{2in}
	\textmd{\textbf{\hmwkClass:\\  \hmwkTitle}}\\
	\normalsize\vspace{0.1in}\small{Due\ on\ \hmwkDueDate\ at 23:59}\\
	\vspace{4in}
}

\author{
	Name: \textbf{\hmwkAuthorName} \\
	Student ID: \hmwkAuthorID}
\date{}

\renewcommand{\part}[1]{\textbf{\large Part \Alph{partCounter}}\stepcounter{partCounter}\\}

%
% Various Helper Commands
%

% Useful for algorithms
\newcommand{\alg}[1]{\textsc{\bfseries \footnotesize #1}}
% For derivatives
\newcommand{\deriv}[1]{\frac{\mathrm{d}}{\mathrm{d}x} (#1)}
% For partial derivatives
\newcommand{\pderiv}[2]{\frac{\partial}{\partial #1} (#2)}
% Integral dx
\newcommand{\dx}{\mathrm{d}x}
% Alias for the Solution section header
\newcommand{\solution}{\textbf{\large Solution}}
% Probability commands: Expectation, Variance, Covariance, Bias
\newcommand{\E}{\mathrm{E}}
\newcommand{\Var}{\mathrm{Var}}
\newcommand{\Cov}{\mathrm{Cov}}
\newcommand{\Bias}{\mathrm{Bias}}

\begin{document}
	
	
	% \maketitle
	% \thispagestyle{empty}
	% \pagebreak
	
	\date{
		Due on Oct. 27, 2026, 18:00 UTC+8}
	\title{SI 140A-01  Probability \& Statistics for EECS, Fall 2026 \\
		Homework 4}
	\maketitle
	Read all the instructions below carefully before you start working on the assignment, and before you make a submission.
	\begin{itemize}
		\item You are required to write down all the major steps towards making your conclusions; otherwise you may obtain limited points of the problem.
		\item Write your homework in English; otherwise you will get no points of this homework.
		\item Any form of plagiarism will lead to $0$ point of this homework. 
	\end{itemize}
	
	
	
	\newpage
  
	\begin{homeworkProblem}
	A book has $n$ typos. Two proofreaders, Prue and Frida, independently read the book.
Prue catches each typo with probability $p_1$ and misses it with probability $q_1 = 1-p_1$, independently, and likewise for Frida, who has probabilities $p_2$ of catching and $q_2=1-p_2$ of missing each typo. Let $X_1$ be the number of typos caught by Prue, $X_2$ be the number caught by Frida, and $X$ be the number caught by at least one of the two proofreaders.
\begin{enumerate}
    \item[(a)] Find the distribution of $X$.
    \item[(b)] For this part only, assume that $p_1 = p_2$. Find the conditional distribution of $X_1$ given that $X_1 + X_2 = t$.
\end{enumerate}


   
		
	\end{homeworkProblem}
	
	\newpage
	\begin{homeworkProblem}
		    Let $X$ be a random day of the week, coded so that Monday is $1$, Tuesday is $2$, \textit{etc}. (so $X$ takes values $1, 2, \dots, 7$ with equal probabilities).
    Let $Y$ be the next day after $X$.
    Do $X$ and $Y$ have the same distribution?
    What is $P(X < Y)$?\\

	\end{homeworkProblem}
	
	
	\newpage
	\begin{homeworkProblem}
	 For $X$ and $Y$ binary digits (0 or 1), let $X \oplus Y$ be 0 if $X=Y$ and 1 if $X\neq Y$ (this operation is called exclusive or (often abbreviated to XOR), or addition mod 2).
\begin{enumerate}[(a)]
    \item  Let $X \sim {\rm Bern}(p)$ and $Y \sim {\rm Bern}(1/2)$, independently. What is the distribution of $X \oplus Y$?
    \item With notation as in sub-problem (a), is $X \oplus Y$ independent of $X$?  Is $X \oplus Y$ independent of $Y$?  Be sure to consider both the case $p= 1/2$ and the case $p\neq 1/2$.
    \item Let $X_1, ... ,X_n$ be i.i.d. (i.e., independent and identically distributed) Bern(1/2) R.V.s. For each nonempty subset $J$ of $\{1,2,...,n\}$, let
    \begin{equation*}
        Y_J=\oplus_{j \in J}X_j.
    \end{equation*}
    
    Show that $Y_J \sim {\rm Bern}(1/2)$ and that these $2^n-1$ R.V.s are pairwise independent, but not independent.
\end{enumerate}
\end{homeworkProblem}


\newpage
\begin{homeworkProblem} 
Let $X$ be a discrete R.V. whose distinct possible values are $a_1, a_2,...,a_n,$ with probabilities $p_1, p_2 ...,p_n,$ respectively (so $p_1 + p_2 + \cdots + p_n = 1$). The entropy of $X$ is defined to be the average surprise of learning the value of $X$:
\begin{equation*}
    H(X)=\sum_{j=1}^np_j\log_2(1/p_j).
\end{equation*}
Show that the maximum entropy for $X$ is when
its distribution is uniform over $a_1, a_2,...,a_n,$ i.e., $p_j = 1/n, \forall j$.
\end{homeworkProblem}

	
\end{document}